\section{AMC 8 1998}
        \begin{enumerate}
        	\item (\textit{AMC 8 1998}) For $x=7$, which of the following is the smallest?


 $\text{(A)}\ \dfrac{6}{x} \qquad \text{(B)}\ \dfrac{6}{x+1} \qquad \text{(C)}\ \dfrac{6}{x-1} \qquad \text{(D)}\ \dfrac{x}{6} \qquad \text{(E)}\ \dfrac{x+1}{6}$

 \textbf{Answer: }B



\textbf{Solution 1}

Plugging $x$ in for every answer choice would give

 $\text{(A)}\ \dfrac{6}{7} \qquad \text{(B)}\ \dfrac{6}{8} \qquad \text{(C)}\ \dfrac{6}{6} \qquad \text{(D)}\ \dfrac{7}{6} \qquad \text{(E)}\ \dfrac{8}{6}$

 From here, we can see that the smallest is answer choice $\boxed{B}$




\textbf{Solution 2}

Note that $\dfrac{6}{x+1}<\dfrac{6}{x}<\dfrac{6}{x-1}$ (for $x>1$) and $\dfrac{x}{6}<\dfrac{x+1}{6}$. Therefore, we just need to compare $\dfrac{6}{x+1}$ and $\dfrac{x}{6}$. Plugging in $x=7$, we get $\dfrac{3}{4}$ and $\dfrac{7}{6}$, respectively, with $\dfrac{3}{4}<\dfrac{7}{6}$. Thus, the answer is $\boxed{(B) \dfrac{6}{x+1}}$. 

 


	\item (\textit{AMC 8 1998}) If $\begin{tabular}{r|l}a&b \\ \hline c&d\end{tabular} = \text{a}\cdot \text{d} - \text{b}\cdot \text{c}$, what is the value of $\begin{tabular}{r|l}3&4 \\ \hline 1&2\end{tabular}$?


 $\text{(A)}\ -2 \qquad \text{(B)}\ -1 \qquad \text{(C)}\ 0 \qquad \text{(D)}\ 1 \qquad \text{(E)}\ 2$

 \textbf{Answer: }E



\textbf{Solution}

Plugging in values for $a$, $b$, $c$, and $d$, we get

 $a=3$, $b=4$, $c=1$, $d=2$, 

 $a\times d=3\times2=6$
$b\times c=4\times1=4$
$6-4=2$
$\boxed{E}$



	\item (\textit{AMC 8 1998}) $\dfrac{\dfrac{3}{8} + \dfrac{7}{8}}{\dfrac{4}{5}} =$
$\text{(A)}\ 1 \qquad \text{(B)} \dfrac{25}{16} \qquad \text{(C)}\ 2 \qquad \text{(D)}\ \dfrac{43}{20} \qquad \text{(E)}\ \dfrac{47}{16}$

 \textbf{Answer: }B



\textbf{Solution}

$\frac{\frac{3}{8}+\frac{7}{8}}{\frac{4}{5}} = \frac{\frac{10}{8}}{\frac{4}{5}} = \frac{\frac{5}{4}}{\frac{4}{5}} = \frac{5}{4} \cdot \frac{5}{4} = \frac{25}{16} = \boxed{B}$



	\item (\textit{AMC 8 1998}) How many triangles are in this figure? (Some triangles may overlap other triangles.)


 
                    \begin{figure}[H]
                        \centering
                        \includegraphics[width=0.5\textwidth]{Problem/AMC_8_AJHSME_1998_4_Problem_Figure_1.png}
                    \end{figure}
                    
$\text{(A)}\ 9 \qquad \text{(B)}\ 8 \qquad \text{(C)}\ 7 \qquad \text{(D)}\ 6 \qquad \text{(E)}\ 5$

 \textbf{Answer: }E



\textbf{Solution}

By inspection, we have that there $5$ triangles:  Each of the $3$ small triangles, $1$ medium triangle made of the rightmost two small triangles, and the $1$ large triangle.

 $\boxed{E}$



	\item (\textit{AMC 8 1998}) Which of the following numbers is largest?


 $\text{(A)}\ 9.12344 \qquad \text{(B)}\ 9.123\overline{4} \qquad \text{(C)}\ 9.12\overline{34} \qquad \text{(D)}\ 9.1\overline{234} \qquad \text{(E)}\ 9.\overline{1234}$

 \textbf{Answer: }B



\textbf{Solution 1}

The first six decimal places of the answer choices are as follows: 

 
            \begin{itemize}
        \item (A): $9.123440$
            \end{itemize}
            
            \begin{itemize}
        \item (B): $9.123444$
            \end{itemize}
            
            \begin{itemize}
        \item (C): $9.123434$
            \end{itemize}
            
            \begin{itemize}
        \item (D): $9.123423$
            \end{itemize}
            
            \begin{itemize}
        \item (E): $9.123412$
            \end{itemize}
            So $\boxed{\textbf{(B)}\ 9.123\overline{4}}$ is the largest of these numbers.

 



\textbf{Solution 2}

Looking at the fifth digit after the decimal of each, we find that choices $C$, $D$, and $E$ are eliminated.Since $B$ has a non-zero sixth digit after the decimal, unlike $A$, we conclude that $\boxed{B}$ is the largest number.

 


	\item (\textit{AMC 8 1998}) Dots are spaced one unit apart, horizontally and vertically.  The number of square units enclosed by the polygon is


 
                    \begin{figure}[H]
                        \centering
                        \includegraphics[width=0.5\textwidth]{Problem/AMC_8_AJHSME_1998_6_Problem_Figure_1.png}
                    \end{figure}
                    
$\text{(A)}\ 5 \qquad \text{(B)}\ 6 \qquad \text{(C)}\ 7 \qquad \text{(D)}\ 8 \qquad \text{(E)}\ 9$

 \textbf{Answer: }B



\textbf{Solution 1}

We could count the area contributed by each square on the $3 \times 3$ grid:

 Top-left: $0$

 Top: Triangle with area $\frac{1}{2}$

 Top-right: $0$

 Left: Square with area $1$

 Center: Square with area $1$

 Right: Square with area $1$

 Bottom-left: Square with area $1$

 Bottom: Triangle with area $\frac{1}{2}$

 Bottom-right: Square with area $1$

 Adding all of these together, we get $6$ which is the same as $\boxed{B}$




\textbf{Solution 2}

By Pick's Theorem, we get the formula, $A=I+\frac{b}{2}-1$ where $I$ is the number of lattice points in the interior and $b$ being the number of lattice points on the boundary. In this problem, we can see that $I=1$ and $B=12$. Substituting gives us $A=1+\frac{12}{2}-1=6$ Thus, the answer is $\boxed{\text{(B) 6}}$




\textbf{Solution 3}

Notice that the extra triangle on the top with area $1/2$ can be placed (like a jigsaw puzzle) at the bottom of the grid where there is a triangular hole, also with area $1/2$. This creates a $2*3$ rectangle, with a area of $6$. The answer is $\boxed{\text{(B) 6}}$~sakshamsethi

 


	\item (\textit{AMC 8 1998}) $100\times 19.98\times 1.998\times 1000=$
$\text{(A)}\ (1.998)^2 \qquad \text{(B)}\ (19.98)^2 \qquad \text{(C)}\ (199.8)^2 \qquad \text{(D)}\ (1998)^2 \qquad \text{(E)}\ (19980)^2$

 \textbf{Answer: }D



\textbf{Solution}

$19.98\times100=1998$
$1.998\times1000=1998$
$1998\times1998=(1998)^2=\boxed{D}$



	\item (\textit{AMC 8 1998}) A child's wading pool contains 200 gallons of water.  If water evaporates at the rate of 0.5 gallons per day and no other water is added or removed, how many gallons of water will be in the pool after 30 days?


 $\text{(A)}\ 140 \qquad \text{(B)}\ 170 \qquad \text{(C)}\ 185 \qquad \text{(D)}\ 198.5 \qquad \text{(E)}\ 199.85$

 \textbf{Answer: }C



\textbf{Solution}

Over the $30$ days, a total of $30 \cdot 0.5 = 15$ gallons of water evaporate. So $200-15=\boxed{\textbf{(C)}\ 185}$ gallons remain. 

 


	\item (\textit{AMC 8 1998}) For a sale, a store owner reduces the price of a $\$10 scarf by $20\%$.  Later the price is lowered again, this time by one-half the reduced price.  The price is now


 $\text{(A)}\ 2.00\text{ dollars} \qquad \text{(B)}\ 3.75\text{ dollars} \qquad \text{(C)}\ 4.00\text{ dollars} \qquad \text{(D)}\ 4.90\text{ dollars} \qquad \text{(E)}\ 6.40\text{ dollars}$

 \textbf{Answer: }C



\textbf{Solution 1}

After the first discount, the price of the scarf is $0.8 \cdot \$10 = \$8$.  

 After the second discount, the price is $\frac{1}{2} \cdot \$8 = \boxed{\textbf{(C)}\ \$4.00}$.

 



\textbf{Solution 2}

The first discount multiplies the price by $1 - 0.2 = 0.8$, while the second discount multiplies the price by $\frac{1}{2}$.  So the two discounts multiply the price by $0.8 \cdot \frac{1}{2} = 0.4$, and thus the final price of the scarf is $0.4 \cdot \$10 = \boxed{\textbf{(C)}\ \$4.00}$.

 


	\item (\textit{AMC 8 1998}) Each of the letters $\text{W}$, $\text{X}$, $\text{Y}$, and $\text{Z}$ represents a different integer in the set $\{ 1,2,3,4\}$, but not necessarily in that order.  If $\dfrac{\text{W}}{\text{X}} - \dfrac{\text{Y}}{\text{Z}}=1$, then the sum of $\text{W}$ and $\text{Y}$ is


 $\text{(A)}\ 3 \qquad \text{(B)}\ 4 \qquad \text{(C)}\ 5 \qquad \text{(D)}\ 6 \qquad \text{(E)}\ 7$

 \textbf{Answer: }E



\textbf{Solution}

There are different ways to approach this problem, and I'll start with the different factor of the numbers of the set $\{ 1,2,3,4\}$.

 $1$ has factor $1$.

 $2$ has factors $1$ and $2$
$3$ has factors $1$ and $3$
$4$ has factors $1$, $2$, and $4$.

 From here, we note that even though all numbers have the factor $1$, only $4$ has another factor other than $1$ in the set (ie. $2$)

 We could therefore have one fraction be $\frac{4}{2}$ and another $\frac{3}{1}$.

 The sum of the numerators is $4+3=7=\boxed{E}$



	\item (\textit{AMC 8 1998}) Harry has 3 sisters and 5 brothers.  His sister Harriet has $\text{S}$ sisters and $\text{B}$ brothers.  What is the product of $\text{S}$ and $\text{B}$?


 $\text{(A)}\ 8 \qquad \text{(B)}\ 10 \qquad \text{(C)}\ 12 \qquad \text{(D)}\ 15 \qquad \text{(E)}\ 18$

 \textbf{Answer: }C



\textbf{Solution}

Since Harry has $3$ sisters and $5$ brothers, his parents have a total of $3$ daughters and $6$ sons.  Therefore, his sister Harriet has $2$ sisters and $6$ brothers, and thus $S \cdot B = 2 \cdot 6 = \boxed{\textbf{(C)}\ 12}$.

 


	\item (\textit{AMC 8 1998}) $2\left(1-\dfrac{1}{2}\right) + 3\left(1-\dfrac{1}{3}\right) + 4\left(1-\dfrac{1}{4}\right) + \cdots + 10\left(1-\dfrac{1}{10}\right)=$
$\text{(A)}\ 45 \qquad \text{(B)}\ 49 \qquad \text{(C)}\ 50 \qquad \text{(D)}\ 54 \qquad \text{(E)}\ 55$

 \textbf{Answer: }A



\textbf{Solution}

Taking the first product, we have

 $\left(1-\frac{1}{2}\right)=\frac{1}{2}$
$\frac{1}{2}\times2=1$

 Looking at the second, we get

 $\left(1-\frac{1}{3}\right)=\frac{2}{3}$
$\frac{2}{3}\times3=2$

 We seem to be going up by $1$.

 Just to check,

 $1-\frac{1}{n}=\frac{n-1}{n}$
$\frac{n-1}{n}\times n=n-1$

 Now that we have discovered the pattern, we have to find the last term.

 $1-\frac{1}{10}=\frac{9}{10}$
$\frac{9}{10}\times10=9$

 The sum of all numbers from $1$ to $9$ is

 $\frac{9\cdot10}{2}=45=\boxed{A}$



	\item (\textit{AMC 8 1998}) What is the ratio of the area of the shaded square to the area of the large square?  (The figure is drawn to scale)


 
                    \begin{figure}[H]
                        \centering
                        \includegraphics[width=0.5\textwidth]{Problem/AMC_8_AJHSME_1998_13_Problem_Figure_1.png}
                    \end{figure}
                    
$\text{(A)}\ \dfrac{1}{6} \qquad \text{(B)}\ \dfrac{1}{7} \qquad \text{(C)}\ \dfrac{1}{8} \qquad \text{(D)}\ \dfrac{1}{12} \qquad \text{(E)}\ \dfrac{1}{16}$

 \textbf{Answer: }C



\textbf{Solution 1}

We can divide the large square into quarters by diagonals.

 Then, in $\frac{1}{4}$ the area of the big square, the little square would have $\frac{1}{2}$ the area.

 $\frac{1}{4}\times\frac{1}{2}=\frac{1}{8}=\boxed{C}$




\textbf{Solution 2}

Answer: \textbf{C}

 
                    \begin{figure}[H]
                        \centering
                        \includegraphics[width=0.5\textwidth]{Solution/AMC_8_AJHSME_1998_13_Solution_Figure_1.png}
                    \end{figure}
                    

 Divide the square into 16 smaller squares as shown. The shaded square is formed from 4 half-squares, so its area is 2. The ratio 2 to 16 is 1/8.

 


	\item (\textit{AMC 8 1998}) At Annville Junior High School, $30\%$ of the students in the Math Club are in the Science Club, and $80\%$ of the students in the Science Club are in the Math Club.  There are 15 students in the Science Club.  How many students are in the Math Club?


 $\text{(A)}\ 12 \qquad \text{(B)}\ 15 \qquad \text{(C)}\ 30 \qquad \text{(D)}\ 36 \qquad \text{(E)}\ 40$

 \textbf{Answer: }E



\textbf{Solution}

If $80\%$ of the people in the science club of $15$ people are in the Math Club, $\frac{4}{5}\times15=12$ people are in the both the Math Club and the Science Club.

 These $12$ people make up $30\%$ of the Math Club.

 Setting up a proportion:

 $12\cdot 1.00 =0.30\cdot x$
$\frac{12}{0.3} = 40=x$

 There are $40=\boxed{E}$ people in the Math Club.

 


	\item (\textit{AMC 8 1998}) \textbf{Don't Crowd the Isles} Problems 15, 16, and 17 all refer to the following:


 
In the very center of the Irenic Sea lie the beautiful Nisos Isles.  In 1998 the number of people on these islands is only 200, but the population triples every 25 years.  Queen Irene has decreed that there must be at least 1.5 square miles for every person living in the Isles.  The total area of the Nisos Isles is 24,900 square miles.


 


 Estimate the population of Nisos in the year 2050.


 $\text{(A)}\ 600 \qquad \text{(B)}\ 800 \qquad \text{(C)}\ 1000 \qquad \text{(D)}\ 2000 \qquad \text{(E)}\ 3000$

 \textbf{Answer: }D



\textbf{Solution}

The population triples every $25$ years from $200$, and there are $50$ years between $2000$ and $2050$, so we will have $200\times3=600\times3=1800$.

 This is an underestimate, since there are actually $2$ more years for the island's population to grow.

 Therefore, $1800$ can be rounded to $\text{(D)}\ 2000$



	\item (\textit{AMC 8 1998}) \textbf{Don't Crowd the Isles} Problems 15, 16, and 17 all refer to the following:


 
In the very center of the Irenic Sea lie the beautiful Nisos Isles.  In 1998 the number of people on these islands is only 200, but the population triples every 25 years.  Queen Irene has decreed that there must be at least 1.5 square miles for every person living in the Isles.  The total area of the Nisos Isles is 24,900 square miles.


 


 Estimate the year in which the population of Nisos will be approximately 6,000.


 $\text{(A)}\ 2050 \qquad \text{(B)}\ 2075 \qquad \text{(C)}\ 2100 \qquad \text{(D)}\ 2125 \qquad \text{(E)}\ 2150$

 \textbf{Answer: }B



\textbf{Solution 1}

We could triple the population every $25$ years and make a chart:

 Year: 2000Population: 200

 Year: 2025Population: 600

 Year: 2050Population: 1800

 Year: 2075Population: 5400

 Year: 2100Population: 16200

 The closest year is 2075, or $\boxed{B}$




\textbf{Solution 2}

We could find out how many periods of 25 years we need to triple by dividing our total from our present number.

 $\frac{6000}{200}=30$

 The power of $3$ that $30$ is closest to is $27=3^3$

 Therefore, after $3$ periods, we will be closest to $6000$.

 $\boxed{B}$



	\item (\textit{AMC 8 1998}) \textbf{Don't Crowd the Isles} Problems 15, 16, and 17 all refer to the following:


 
In the very center of the Irenic Sea lie the beautiful Nisos Isles.  In 1998 the number of people on these islands is only 200, but the population triples every 25 years.  Queen Irene has decreed that there must be at least 1.5 square miles for every person living in the Isles.  The total area of the Nisos Isles is 24,900 square miles.


 


 In how many years, approximately, from 1998 will the population of Nisos be as much as Queen Irene has proclaimed that the islands can support?


 $\text{(A)}\ 50\text{ yrs.} \qquad \text{(B)}\ 75\text{ yrs.} \qquad \text{(C)}\ 100\text{ yrs.} \qquad \text{(D)}\ 125\text{ yrs.} \qquad \text{(E)}\ 150\text{ yrs.}$

 \textbf{Answer: }C



\textbf{Solution 1}

We can divide the total area by how much will be occupied per person:

 $\frac{24900 \text{ acres}}{1.5 \text{ acres per person}}=16600 \text{ people}$ can stay on the island at its maximum capacity.

 We can divide 16600 by the current population in $1998$ which is 200 to see by what factor the population increases:

 $\frac{16600}{200}=83$-fold increase in population.

 Thus, the population increases by a factor $83$.  This is very close to $3 \times 3 \times 3 \times 3 = 81$, and so there are about $4$ triplings of the island's population.

 It takes $4\times25=100=\boxed{C}$ years to triple the island's population four times in succession.

 



\textbf{Solution 2}

We can continue the pattern, and because the pattern increases numbers rapidly, it won't be hard. $\frac{24900 \text{ acres}}{1.5 \text{ acres per person}}=16600 \text{ people}$ can live on the island at its maximum capacity. 

 $1998: 200 \text{ people}$
$2023: 600 \text{ people}$
$2048: 1800 \text{ people}$
$2073: 5400 \text{ people}$
$2098: 16200 \text{ people}$
After the year $2098$, it will not be possible for the next increase to occur, because when $16200$ is tripled, it is way more than the maximum capacity.
Thus, the answer is $2098-1998 = 100$, or $\boxed{C}$



	\item (\textit{AMC 8 1998}) As indicated by the diagram below, a rectangular piece of paper is folded bottom to top, then left to right, and finally, a hole is punched at X.  What does the paper look like when unfolded?


 
                    \begin{figure}[H]
                        \centering
                        \includegraphics[width=0.5\textwidth]{Problem/AMC_8_AJHSME_1998_18_Problem_Figure_1.png}
                    \end{figure}
                    

                    \begin{figure}[H]
                        \centering
                        \includegraphics[width=0.5\textwidth]{Problem/AMC_8_AJHSME_1998_18_Problem_Figure_2.png}
                    \end{figure}
                    

 \textbf{Answer: }B



\textbf{Solution 1}

By reversing the folds, we find that the holes will be

 1)  In the middle (in relation from left to right)

 2)  At separating poles (up and down)

 This is pictured in $\boxed{B}$




\textbf{Solution 2}

If you unfold the paper, the hole you punched will not move.  The only answer that keeps the hole in the upper-right quadrant in the same place is $\boxed{B}$.

 



\textbf{Solution 3}

You can check each answer choice by folding the paper the same way it was folded before the hole was punched. The only answer choice that has the hole in the right spot is $\boxed{B}$.

 


	\item (\textit{AMC 8 1998}) Tamika selects two different numbers at random from the set $\{ 8,9,10 \}$ and adds them.  Carlos takes two different numbers at random from the set $\{3, 5, 6\}$ and multiplies them.  What is the probability that Tamika's result is greater than Carlos' result?


 $\text{(A)}\ \dfrac{4}{9} \qquad \text{(B)}\ \dfrac{5}{9} \qquad \text{(C)}\ \dfrac{1}{2} \qquad \text{(D)}\ \dfrac{1}{3} \qquad \text{(E)}\ \dfrac{2}{3}$

 \textbf{Answer: }A



\textbf{Solution}

The different possible (and equally likely) values Tamika gets are:

 $8+9=17$
$8+10=18$
$9+10=19$

 The different possible (and equally likely) values Carlos gets are:

 $3\times5=15$
$3\times6=18$
$5\times6=30$

 The probability that if Tamika had the sum $17$ her sum would be greater than Carlos's set is $\frac{1}{3}$, because $17$ is only greater than $15$

 The probability that if Tamika had the sum $18$ her sum would be greater than Carlos's set is $\frac{1}{3}$, because $18$ is only greater than $15$

 The probability that if Tamika had the sum $19$ her sum would be greater than Carlos's set is $\frac{2}{3}$, because $19$ is greater than both $15$ and $18$

 Each sum has a $\frac{1}{3}$ possibility of being chosen, so we have

 $\frac{1}{3}\left(\frac{1}{3}+\frac{1}{3}+\frac{2}{3}\right)=\frac{1}{3}\left(\frac{4}{3}\right)=\frac{4}{9}=\boxed{A}$



	\item (\textit{AMC 8 1998}) Let $PQRS$ be a square piece of paper.  $P$ is folded onto $R$ and then $Q$ is folded onto $S$.  The area of the resulting figure is 9 square inches.  Find the perimeter of square $PQRS$.


 
                    \begin{figure}[H]
                        \centering
                        \includegraphics[width=0.5\textwidth]{Problem/AMC_8_AJHSME_1998_20_Problem_Figure_1.png}
                    \end{figure}
                    
$\text{(A)}\ 9 \qquad \text{(B)}\ 16 \qquad \text{(C)}\ 18 \qquad \text{(D)}\ 24 \qquad \text{(E)}\ 36$

 \textbf{Answer: }D



\textbf{Solution}

After both folds are completed, the square would become a triangle that has an area of $\frac{1}{2} \cdot \frac{1}{2} = \frac{1}{4}$ of the original square.

 Since the area is $9$ square inches for $\frac{1}{4}$ of the square, $9\times4=36$ square inches is the area of square $PQRS$

 The length of the side of a square that has an area of $36$ square inches is $\sqrt{36}=6$ inches.

 Each side is $6$ inches, so the total perimeter is $6\times4=24=\boxed{D}$



	\item (\textit{AMC 8 1998}) A $4\times 4\times 4$ cubical box contains 64 identical small cubes that exactly fill the box.  How many of these small cubes touch a side or the bottom of the box?


 $\text{(A)}\ 48 \qquad \text{(B)}\ 52 \qquad \text{(C)}\ 60 \qquad \text{(D)}\ 64 \qquad \text{(E)}\ 80$

 \textbf{Answer: }B



\textbf{Solution 1}

The box is filled with four $4 \times 4$ layers of cubes.

 All $16$ cubes in the bottom layer touch the bottom of the box.

 In each of the remaining $3$ layers, all the cubes except for those in the center $2 \times 2$ square touch the side of the box.  Thus $12$ cubes in each of these layers touch the box.

 Therefore, a total of $16 + 3 \cdot 12 = \boxed{\textbf{(B)}\ 52}$ cubes touch the box.

 



\textbf{Solution 2}

The box has a total of $4^3 = 64$ cubes, arranged in four $4 \times 4$ layers.

 The only cubes which do not touch the box are those in the middle $2 \times 2$ squares of each of the top $3$ layers.  So $2 \cdot 2 \cdot 3 = 12$ cubes do not touch the box, and thus $64 - 12 = \boxed{\textbf{(B)}\ 52}$ cubes touch the box.

 



\textbf{Solution 3}

Each small cube would have dimensions $1\times 1\times 1$ making each cube a unit cube.

 If there are $16$ cubes per face and there are $5$ faces we are counting, we have $16\times 5= 80$ cubes.

 Some cubes are on account of overlap between different faces.

 We could reduce this number by subtracting the overlap areas, which could mean subtracting 4 cubes from each side and 12 from the bottom.

 $80-(4\times 4+12)=80-28=52=\boxed{B}$




\textbf{Solution 4}

We can use casework to solve this.

 Case $1$: The cubes on the sides

 There are four sides of a cube which means $64$ cubes, but there will be an overlap of $16$ cubes ($4$ sides $\times$ $4$ cubes on each side).So there are $16 \times 4-16=64-16=48$ cubes on the sides.

 Case $2$: The cubes on the bottom

 There are $16$ cubes on the bottom, but by counting the cubes on the sides, $12$ of the cubes are already accounted for. $16-12=4$ cubes on the bottom that aren't already counted. 

 $48+4=\boxed{\textbf{(B)} 52}$ cubes touching at least one side. Since there is no top, we don't need to count the top face cubes.

 


	\item (\textit{AMC 8 1998}) Terri produces a sequence of positive integers by following three rules.  She starts with a positive integer, then applies the appropriate rule to the result, and continues in this fashion.  


 



 Rule 1: If the integer is less than 10, multiply it by 9.


 Rule 2: If the integer is even and greater than 9, divide it by 2.


 Rule 3: If the integer is odd and greater than 9, subtract 5 from it.


 



 A sample sequence: $23, 18, 9, 81, 76, \ldots .$


 Find the $98^\text{th}$ term of the sequence that begins $98, 49, \ldots .$
$\text{(A)}\ 6 \qquad \text{(B)}\ 11 \qquad \text{(C)}\ 22 \qquad \text{(D)}\ 27 \qquad \text{(E)}\ 54$

 \textbf{Answer: }D



\textbf{Solution}

Since each term after the first depends solely on the previous term, the sequence will repeat the same subsequence endlessly once a term appears twice.

 The first terms of the sequence are $98, 49, 44, 22, 11, 6, 54, 27, 22$, and thus the $5$-term subsequence $22, 11, 6, 54, 27$ repeats endlessly starting at the $4^{th}$ term.  Since the $4^{th}$ term is $22$, the $94^{th}$ term is also $22$, and thus the $98^{th}$ term is $\boxed{\textbf{(D)}\ 27}$.

 


	\item (\textit{AMC 8 1998}) If the pattern in the diagram continues, what fraction of the interior would be shaded in the eighth triangle?


 
                    \begin{figure}[H]
                        \centering
                        \includegraphics[width=0.5\textwidth]{Problem/AMC_8_AJHSME_1998_23_Problem_Figure_1.png}
                    \end{figure}
                    
$\text{(A)}\ \dfrac{3}{8} \qquad \text{(B)}\ \dfrac{5}{27} \qquad \text{(C)}\ \dfrac{7}{16} \qquad \text{(D)}\ \dfrac{9}{16} \qquad \text{(E)}\ \dfrac{11}{45}$

 \textbf{Answer: }C



\textbf{Solution}

Notice that the $n^{th}$ triangle contains $n$ rows of congruent smaller triangles, and for each integer $k$ from $1$ to $n$ inclusive, the $k^{th}$ row contains $k$ unshaded and $k-1$ shaded small triangles.

 Therefore, the $8^{th}$ triangle contains $8$ rows, which contain a total of $1 + 2 + 3 + ... + 8 = 36$ unshaded small triangles and $0 + 1 + 2 + ... + 7 = 28$ shaded small triangles.  Therefore, $\frac{28}{28 + 36} = \frac{28}{64} = \boxed{\textbf{(C)}\ \frac{7}{16}}$ of the $8^{th}$ triangle is shaded.

 


	\item (\textit{AMC 8 1998}) A rectangular board of 8 columns has squares numbered beginning in the upper left corner and moving left to right so row one is numbered 1 through 8, row two is 9 through 16, and so on.  A student shades square 1, then skips one square and shades square 3, skips two squares and shades square 6, skips 3 squares and shades square 10, and continues in this way until there is at least one shaded square in each column.  What is the number of the shaded square that first achieves this result?


 
                    \begin{figure}[H]
                        \centering
                        \includegraphics[width=0.5\textwidth]{Problem/AMC_8_AJHSME_1998_24_Problem_Figure_1.png}
                    \end{figure}
                    
$\text{(A)}\ 36 \qquad \text{(B)}\ 64 \qquad \text{(C)}\ 78 \qquad \text{(D)}\ 91 \qquad \text{(E)}\ 120$

 \textbf{Answer: }E



\textbf{Solution 1}

The numbers that are shaded are the triangular numbers, which are numbers in the form $\frac{(n)(n+1)}{2}$ for positive integers.  They can also be generated by starting with $1$, and adding $1, 2, 3, 4...$ as in the description of the problem.

 Squares that have the same remainder after being divided by $8$ will be in the same column.  Thus, we want to find when the last remainder, from $0$ to $7$, is found.

 So, instead of adding $1, 2, 3, 4, 5, 6, 7, 8, 9, 10...$, we can effectively either add $1, 2, 3, 4, 5, 6, 7, 0, 1, 2, 3...$ or subtract $7, 6, 5, 4, 3, 2, 1, 0, 7, 6, 5...$ if we are only concerned about remainders when divided by $8$.  We will pick the number that keeps the terms on the list between $1$ and $8$.  We get:

 $1$
$1 + 2 = 3$
$3 + 3 = 6$
$6 - 4 = 2$
$2 + 5 = 7$
$7 - 2 = 5$
$5 - 1 = 4$
$4 + 0 = 4$
$4 + 1 = 5$
$5 + 2 = 7$
$7 - 5 = 2$
$2 + 4 = 6$
$6 - 3 = 3$
$3 - 2 = 1$
$1 - 1 = 0$

 Finally, a term with $0$ is found, and checking, all numbers $1$ through $7$ are also on the right side of the list.  This means the last term in our sequence is the first time that column $8$ is shaded.  There are $15$ terms in the sequence, leading to an answer of $\frac{15\cdot 16}{2} = 120$, which is choice $\boxed{E}$.

 



\textbf{Solution 2}

Note that the triangular numbers up to $120$ are $1, 3, 6, 10, 15, 21, 28, 36, 45, 55, 66, 78, 91, 105, 120$.  When you divide each of those numbers by $8$, all remainders must be present.  We first search for number(s) that are evenly divisible by $8$; if two such numbers exist, we search for numbers that leave a remainder of $1$, etc.

 Quickly scanning the list, only $6, 10, 28, 36, 66, 78$ and $120$ are even.  That smaller list doesn't have any multiples of $8$ until it hits $120$.  So $\boxed{\textbf{(E)}\ 120}$ must be the answer.

 



\textbf{Solution 3}

The student shades cells containing triangular numbers in increasing order starting from $1$.  Thus the first cells the student shades contain $1$, $3$, $6$, $10$, $15$, $21$, and $28$, so they are in the first, third, sixth, second, seventh, fifth, and fourth columns respectively.  

 So after shading $7$ cells, only the eighth column, whose cells contain multiples of $8$, lacks a shaded cell.  Therefore, the last cell the student shades contains the smallest triangular number divisible by 8.

 Recall that the integer $n$ is triangular if and only if $n = \frac{k(k+1)}{2}$ for some positive integer $k$. So if $n$ is divisible by $8$, $2n = k(k+1)$ is divisible by $16$.  Since only one of $k$ and $k+1$ can be even, either $k$ or $k+1$ must be divisible by $16$. The smallest such integer $k$ such this holds is $15$, and thus the smallest triangular number which is divisible by $8$ is  $\boxed{\textbf{(E)}\ 120}$.

 



\textbf{Solution 4 (Answer Choices)}

The squares the student shades contain the triangular numbers, so the sixth shaded square is $21$, which is in the fifth column, while the seventh shaded square is $28$, which is in the fourth column.  So after the student has shaded $7$ squares, only the eighth column, which contains all positive integers which are divisible by $8$, has no unshaded squares.  Therefore, the last square the student shades must contain a number divisible by $8$.  

 Of the answer choices, only $64$ and $120$ are divisible by $8$. Since $64$ is not a triangular number, the correct answer is $\boxed{\textbf{(E)}\ 120}$.

 


	\item (\textit{AMC 8 1998}) Three generous friends, each with some cash, redistribute their money as follows: Amy gives enough money to Jan and Toy to double the amount that each has.  Jan then gives enough to Amy and Toy to double their amounts.  Finally, Toy gives Amy and Jan enough to double their amounts.  If Toy has $36 when they begin and $36 when they end, what is the total amount that all three friends have?


 $\text{(A)}\ 108\text{ dollars} \qquad \text{(B)}\ 180\text{ dollars} \qquad \text{(C)}\ 216\text{ dollars} \qquad \text{(D)}\ 252\text{ dollars} \qquad \text{(E)}\ 288\text{ dollars}$

 \textbf{Answer: }D



\textbf{Solution}

After Ami gives money to Toy and Jan, Toy has $2 \cdot \$36 = \$72$.

 After Jan gives money to Ami and Toy, Toy has $2 \cdot \$72 = \$144$.

 Therefore, Toy gives $\$144 - \$36 = \$108$ to Jan and Ami.  Since Jan and Ami each doubled their money, they must have had a total of $\$108$ before Toy gave them money and thus now have a total of $2 \cdot \$108 = \$216$.  

 Therefore, the three friends have a total of $\$216 + \$36 = \boxed{\textbf{(D)}\ \$252}$.

 


\end{enumerate}